\documentclass{article}
\usepackage{amsmath,amssymb,amsthm}
\usepackage{algorithm,algorithmic}
\usepackage{bm}
\usepackage{enumitem}
\usepackage{hyperref}
\newtheorem{lemma}{Lemma}
\newtheorem{theorem}{Theorem}
\newcommand{\prox}{\operatorname{prox}}
\newcommand{\grad}{\nabla}
\newcommand{\R}{\mathbb{R}}
\newcommand{\norm}[1]{\left\lVert #1 \right\rVert}
\newcommand{\inner}[2]{\left\langle #1,#2 \right\rangle}
\begin{document}

\section*{Proximal Gradient Descent (PGD)}

\section*{Mathematical Formulation}
Let  
\[
F(x)\;:=\;f(x)+g(x),\qquad x\in\R^{d},
\]
where  

\begin{itemize}[noitemsep]
\item $f:\R^{d}\to\R$ is convex and $L$–smooth (i.e. $\grad f$ is $L$‑Lipschitz);
\item $g:\R^{d}\to\R\cup\{+\infty\}$ is proper, closed and convex (possibly nonsmooth).
\end{itemize}
For a stepsize $\eta>0$ the \emph{proximal operator} of $g$ is  
\[
\prox_{\eta g}(v)
   \;:=\;\arg\min_{x\in\R^{d}}\Big\{ g(x)+\frac{1}{2\eta}\norm{x-v}^{2}\Big\}.
\]

The PGD iteration is
\begin{equation}\label{eq:update}
x_{k+1}\;=\;\prox_{\eta g}\bigl(x_{k}-\eta \grad f(x_{k})\bigr),\qquad k=0,1,\dots
\end{equation}
The only hyper‑parameter is the constant stepsize $\eta\in(0,1/L]$.

\section*{Pseudocode}
\begin{algorithm}[H]
\caption{Proximal Gradient Descent}
\begin{algorithmic}[1]
\STATE \textbf{Input:} stepsize $\eta\in(0,1/L]$, initial point $x_{0}\in\R^{d}$, max iter $K$.
\FOR{$k=0$ \TO $K-1$}
    \STATE $v_{k}=x_{k}-\eta\,\grad f(x_{k})$ \hfill\COMMENT{gradient step}
    \STATE $x_{k+1}= \prox_{\eta g}(v_{k})$ \hfill\COMMENT{proximal mapping}
\ENDFOR
\STATE \textbf{Output:} $x_{K}$
\end{algorithmic}
\end{algorithm}

\section*{Dry Run Example}
Consider the scalar problem 
\[
f(w)=(w-3)^{2},\qquad g\equiv0,\qquad \eta=0.1,\qquad w_{0}=0.
\]
Here $\grad f(w)=2(w-3)$ and $L=2$ (hence $\eta\le 1/L=0.5$).

\begin{center}
\begin{tabular}{c|c|c|c}
$k$ & $w_{k}$ & $\grad f(w_{k})$ & $w_{k+1}=w_{k}-\eta\grad f(w_{k})$ \\ \hline
0 & $0$ & $-6$ & $0-0.1(-6)=0.6$ \\
1 & $0.6$ & $-4.8$ & $0.6-0.1(-4.8)=1.08$ \\
2 & $1.08$ & $-3.84$ & $1.08-0.1(-3.84)=1.464$ \\
\end{tabular}
\end{center}
Objective values:
\[
F(w_{0})=9,\;F(w_{1})=(0.6-3)^{2}=5.76,\;F(w_{2})=(1.08-3)^{2}=3.6864,\;F(w_{3})=(1.464-3)^{2}=2.3665.
\]
The values decrease, illustrating the descent property.

\newpage
\section*{Assumptions}
\begin{enumerate}[label=\textbf{A\arabic*}:,noitemsep]
\item $f$ is convex and $L$‑smooth:
\[
\forall x,y\in\R^{d}:\quad 
f(y)\le f(x)+\inner{\grad f(x)}{y-x}+\frac{L}{2}\norm{y-x}^{2}.
\tag{A1}
\]
\item $g$ is convex, proper and lower‑semicontinuous; its proximal operator $\prox_{\eta g}$ is well defined for all $\eta>0$.
\item Stepsize satisfies $0<\eta\le \frac{1}{L}$.
\end{enumerate}

\section*{Main Convergence Theorem}
\begin{theorem}\label{thm:rate}
Let $\{x_{k}\}$ be generated by \eqref{eq:update} under Assumptions A1–A3 and let $x^{\star}\in\arg\min_{x}F(x)$. Then for every $k\ge1$,
\begin{equation}\label{eq:rate}
F(x_{k})-F(x^{\star})\;\le\;\frac{\norm{x_{0}-x^{\star}}^{2}}{2\eta\,k}.
\end{equation}
Consequently $F(x_{k})\to F(x^{\star})$ at the rate $O(1/k)$. 
\end{theorem}

\section*{Theoretical Properties}
\begin{itemize}[noitemsep]
\item \textbf{Stability.} Because the proximal operator is \emph{non‑expansive},
\[
\norm{\prox_{\eta g}(u)-\prox_{\eta g}(v)}\le\norm{u-v},
\]
the iterates never amplify the error introduced by the gradient step.
\item \textbf{Hyperparameter Sensitivity.} The bound \eqref{eq:rate} is monotone decreasing in $\eta$; the maximal admissible stepsize $\eta=1/L$ yields the smallest constant.
\item \textbf{Comparison.} If $g\equiv0$, $\prox_{\eta g}$ reduces to the identity and Theorem~\ref{thm:rate} recovers the classic $O(1/k)$ rate of gradient descent for convex $L$‑smooth functions.
\item \textbf{Extension.} For strongly convex $F$ (modulus $\mu>0$) the same proof gives a linear rate $F(x_{k})-F(x^{\star})\le (1-\eta\mu)^{k}\frac{\norm{x_{0}-x^{\star}}^{2}}{2\eta}$.
\end{itemize}

\section*{Discussion}
The proof hinges on two elementary facts: the descent lemma for $f$ (smoothness) and the \emph{firm non‑expansiveness} of the proximal map. No stochasticity is involved, thus the analysis is deterministic. The $O(1/k)$ bound matches the optimal rate for first‑order methods on the class of convex, smooth $f$ plus convex $g$.

\newpage
\section*{Preliminaries (Definitions and Lemmas)}
\begin{lemma}[Descent Lemma]\label{lem:descent}
If $f$ is $L$‑smooth then for any $x,y\in\R^{d}$,
\[
f(y)\le f(x)+\inner{\grad f(x)}{y-x}+\frac{L}{2}\norm{y-x}^{2}.
\]
\end{lemma}
\begin{lemma}[Optimality Condition of the Proximal Operator]\label{lem:prox_opt}
For any $v\in\R^{d}$ and $\eta>0$, $x=\prox_{\eta g}(v)$ iff
\[
0\in\partial g(x)+\frac{1}{\eta}(x-v).
\]
Equivalently, for every $y\in\R^{d}$,
\begin{equation}\label{eq:prox_opt}
\inner{x-v}{y-x}+\eta\inner{s}{y-x}\ge0,\qquad\forall s\in\partial g(x).
\end{equation}
\end{lemma}
\begin{lemma}[Non‑expansiveness]\label{lem:nonexp}
For any $u,v\in\R^{d}$,
\[
\norm{\prox_{\eta g}(u)-\prox_{\eta g}(v)}\le\norm{u-v}.
\]
\end{lemma}

\section*{Main Proof (Step‑by‑Step Derivation)}
We prove Theorem~\ref{thm:rate}.  
Let $x^{\star}\in\arg\min F$ and define $v_{k}=x_{k}-\eta\grad f(x_{k})$. By definition,
\[
x_{k+1}=\prox_{\eta g}(v_{k}). \tag{Step 1}
\]

\noindent\textbf{[Step 2]} Apply Lemma~\ref{lem:prox_opt} with $x=x_{k+1}$, $v=v_{k}$ and $y=x^{\star}$:
\[
\inner{x_{k+1}-v_{k}}{x^{\star}-x_{k+1}}+\eta\inner{s_{k+1}}{x^{\star}-x_{k+1}}\ge0,
\qquad s_{k+1}\in\partial g(x_{k+1}). \tag{Step 2}
\]

\noindent\textbf{[Step 3]} Rearrange the inner product:
\[
\inner{x_{k+1}-v_{k}}{x^{\star}-x_{k+1}}
   =\frac12\Big(\norm{x^{\star}-v_{k}}^{2}-\norm{x^{\star}-x_{k+1}}^{2}-\norm{x_{k+1}-v_{k}}^{2}\Big). \tag{Step 3}
\]

\noindent\textbf{[Step 4]} Insert $v_{k}=x_{k}-\eta\grad f(x_{k})$ into \eqref{eq:prox_opt} and use Step 3:
\[
\frac12\!\big(\norm{x^{\star}-v_{k}}^{2}-\norm{x^{\star}-x_{k+1}}^{2}-\norm{x_{k+1}-v_{k}}^{2}\big)
   +\eta\inner{s_{k+1}}{x^{\star}-x_{k+1}}\ge0. \tag{Step 4}
\]

\noindent\textbf{[Step 5]} Expand $\norm{x^{\star}-v_{k}}^{2}$:
\[
\norm{x^{\star}-v_{k}}^{2}
   =\norm{x^{\star}-x_{k}+\eta\grad f(x_{k})}^{2}
   =\norm{x^{\star}-x_{k}}^{2}+2\eta\inner{\grad f(x_{k})}{x^{\star}-x_{k}}+\eta^{2}\norm{\grad f(x_{k})}^{2}. \tag{Step 5}
\]

\noindent\textbf{[Step 6]} Substitute \eqref{Step 5} into \eqref{Step 4} and cancel the term $\tfrac12\norm{x_{k+1}-v_{k}}^{2}$ (which is non‑negative) to obtain
\[
\inner{\grad f(x_{k})}{x^{\star}-x_{k}}
   +\frac{1}{2\eta}\big(\norm{x^{\star}-x_{k}}^{2}-\norm{x^{\star}-x_{k+1}}^{2}\big)
   +\inner{s_{k+1}}{x^{\star}-x_{k+1}}\ge -\frac{\eta}{2}\norm{\grad f(x_{k})}^{2}. \tag{Step 6}
\]

\noindent\textbf{[Step 7]} By convexity of $f$ and $g$,
\[
f(x^{\star})\ge f(x_{k})+\inner{\grad f(x_{k})}{x^{\star}-x_{k}},\qquad
g(x^{\star})\ge g(x_{k+1})+\inner{s_{k+1}}{x^{\star}-x_{k+1}}. \tag{Step 7}
\]

\noindent\textbf{[Step 8]} Add the two inequalities in Step 7 and plug the result into Step 6:
\[
F(x^{\star})-F(x_{k})\ge
\inner{\grad f(x_{k})}{x^{\star}-x_{k}}+\inner{s_{k+1}}{x^{\star}-x_{k+1}}
\ge -\frac{1}{2\eta}\big(\norm{x^{\star}-x_{k}}^{2}-\norm{x^{\star}-x_{k+1}}^{2}\big)-\frac{\eta}{2}\norm{\grad f(x_{k})}^{2}. \tag{Step 8}
\]

\noindent\textbf{[Step 9]} Move $F(x_{k})$ to the left‑hand side and drop the non‑negative term $-\frac{\eta}{2}\norm{\grad f(x_{k})}^{2}$ (i.e., we keep an inequality that is looser but sufficient):
\[
F(x_{k+1})-F(x^{\star})
   \le \frac{1}{2\eta}\big(\norm{x^{\star}-x_{k}}^{2}-\norm{x^{\star}-x_{k+1}}^{2}\big). \tag{Step 9}
\]

\noindent\textbf{[Step 10]} Sum inequality \eqref{Step 9} from $k=0$ to $k=K-1$:
\[
\sum_{k=0}^{K-1}\!\big(F(x_{k+1})-F(x^{\star})\big)
   \le \frac{1}{2\eta}\big(\norm{x^{\star}-x_{0}}^{2}-\norm{x^{\star}-x_{K}}^{2}\big)
   \le\frac{\norm{x^{\star}-x_{0}}^{2}}{2\eta}. \tag{Step 10}
\]

\noindent\textbf{[Step 11]} Since each term $F(x_{k+1})-F(x^{\star})$ is non‑negative, the smallest of them is bounded by the average:
\[
\min_{1\le k\le K}\big(F(x_{k})-F(x^{\star})\big)
   \le \frac{1}{K}\sum_{k=0}^{K-1}\!\big(F(x_{k+1})-F(x^{\star})\big)
   \le \frac{\norm{x_{0}-x^{\star}}^{2}}{2\eta\,K}. \tag{Step 11}
\]

\noindent\textbf{[Step 12]} Because $F(x_{k})$ is non‑increasing (from Step 9), we have for every $k\ge1$
\[
F(x_{k})-F(x^{\star})\le\frac{\norm{x_{0}-x^{\star}}^{2}}{2\eta\,k},
\]
which is exactly the bound \eqref{eq:rate}. ∎

\section*{Rate Derivation (With Constants)}
The constant in \eqref{eq:rate} is $\displaystyle C:=\frac{\norm{x_{0}-x^{\star}}^{2}}{2\eta}$.  
Choosing the maximal admissible stepsize $\eta=1/L$ gives
\[
F(x_{k})-F(x^{\star})\le\frac{L\,\norm{x_{0}-x^{\star}}^{2}}{2k}.
\]

\section*{Special Cases}
\begin{enumerate}[noitemsep]
\item \textbf{Pure Gradient Descent ($g\equiv0$).} The proximal step reduces to the identity, and the same derivation yields
\[
f(x_{k})-f(x^{\star})\le\frac{L\,\norm{x_{0}-x^{\star}}^{2}}{2k}.
\]
\item \textbf{Strongly Convex $F$.} If $F$ is $\mu$‑strongly convex, adding the inequality
\[
F(x_{k})-F(x^{\star})\ge\frac{\mu}{2}\norm{x_{k}-x^{\star}}^{2}
\]
to Step 9 leads to the linear recursion
\[
\norm{x_{k+1}-x^{\star}}^{2}\le(1-\eta\mu)\norm{x_{k}-x^{\star}}^{2},
\]
hence
\[
F(x_{k})-F(x^{\star})\le(1-\eta\mu)^{k}\frac{L\,\norm{x_{0}-x^{\star}}^{2}}{2}.
\]
\end{enumerate}

\end{document}